\clearpage
\pagenumbering{harfi}
\setcounter{page}{1}
\pagestyle{plain}

% --- صفحه تقدیم به ---
\mbox{}\\
\noindent{\fontsize{16}{22}\selectfont\textbf{تقدیم به}}
\addcontentsline{toc}{section}{تقدیم به}

% \vspace{2.5cm}

% \begin{center}
%     {\fontsize{16}{24}\selectfont\textbf{پدر و مادر عزیزم؛}}\\[0.6cm]
%     {\fontsize{14}{22}\selectfont
%     که با حمایت‌های بی‌دریغ، صبر و شکیبایی، همواره راهنما و پشتوانهٔ من در تمام مراحل زندگی و تحصیل بوده‌اند.
%     }
% \end{center}

% --- صفحه تشکر و قدردانی ---
\clearpage
\mbox{}\\
\noindent{\fontsize{16}{22}\selectfont\textbf{تشکر و قدردانی}}
\addcontentsline{toc}{section}{تشکر و قدردانی}

% \vspace{1.2cm}
% \noindent
% بر خود لازم می‌دانم از استاد راهنمای ارجمندم، \textbf{جناب آقای دکتر صدیق}، که با راهنمایی‌های ارزشمند، صبر و شکیبایی، و حمایت‌های بی‌دریغ خود در تمامی مراحل انجام این پروژه راهنما و پشتیبان من بودند، کمال تشکر و قدردانی را داشته باشم.

% \vspace{0.4cm}
% \noindent
% همچنین از استاد گرامی و فرهیخته، \textbf{جناب آقای دکتر شریعت‌پناهی}، به پاس رهنمودها و دانش ارزشمندی که در طول دوران تحصیل در اختیار اینجانب قرار دادند، صمیمانه سپاسگزارم.

% \vspace{0.6cm}
% \noindent
% برای این اساتید بزرگوار از درگاه خداوند متعال سلامتی و توفیق روزافزون آرزومندم.

% --- صفحه چکیده ---
\clearpage
\mbox{}\\% معادل یک خط خالی در اول صفحه مثل ورد
\noindent{\fontsize{16}{22}\selectfont\textbf{چکیده}}
\addcontentsline{toc}{section}{چکیده}

\vspace{0.8cm}
\noindent
هر گزارش پروژه یا هر پایان‌نامه‌ای با چکیده آغاز می‌شود که عبارت است از توضیح راجع به موضوع پروژه، شیوه‌ها و روش انجام پروژه و نتیجه نهایی. کلمات یا عباراتی که در این بخش توضیح داده می‌شود، باید کاملاً محوری و مرتبط با موضوع پروژه باشند. در متن چکیده، از ارجاع به منابع و اشاره به جداول و نمودارها اجتناب شود.

\vspace{1cm}
\noindent\textbf{واژه‌های کلیدی:} تعداد کلمات یا عبارات کلیدی حداکثر می‌تواند پنج کلمه یا عبارت باشد.

% --- ۱. صفحه فهرست مطالب ---
\clearpage
\mbox{}\\% معادل یک خط خالی در اول صفحه مثل ورد
\tableofcontents

% --- ۲. صفحه فهرست اشکال ---
\clearpage
\mbox{}\\% معادل یک خط خالی در اول صفحه مثل ورد
\listoffigures

% --- ۳. صفحه فهرست جداول ---
\clearpage
\mbox{}\\% معادل یک خط خالی در اول صفحه مثل ورد
\listoftables

% --- ۴. صفحه فهرست علائم اختصاری ---
\clearpage
\mbox{}\\% معادل یک خط خالی در اول صفحه مثل ورد
\noindent{\fontsize{16}{22}\selectfont\textbf{فهرست علائم اختصاری}}
\addcontentsline{toc}{section}{فهرست علائم اختصاری}
\vspace{0.8cm}

\begin{table}[h!]
    \centering
    \renewcommand{\arraystretch}{1.5}
    \begin{tabular}{l@{\hspace{3cm}}r}
        \toprule
        \textbf{کمیت و واحد} & \textbf{شرح فارسی} \\
        \midrule
        $v \ (\text{m/s})$   & سرعت               \\
        $a \ (\text{m/s}^2)$ & شتاب               \\
        $F \ (\text{N})$     & نیرو               \\
        \bottomrule
    \end{tabular}
\end{table}